\zlabel{triptych:concise:commentary:start}
\section{What the servants can do while they wait}

\subsection{How one people can speak in several voices}

If we took the loudest voice in Corinth for the voice of the entire church, Paul's thanksgiving would make little sense. He says the Corinthians lack no gift. Yet within a few lines they are quarreling, and later in the letter he will have to speak to some as to children (1 Cor 1:7, 10--13 and 3:1). Something like this difficulty occurred to Jerome as he read Isaiah. Both a humble appeal to the Father and what sounds like a complaint about his conduct are made in the first person plural. Who is entitled to say what God has done to all of them?

Jerome turns to the opening of Paul's letter. The address includes those who call on the Lord wherever they are. In the body of the letter Paul can then speak quite particularly to the unmarried or to people in danger of idolatry. The whole church is entitled to hear what is addressed to each. An assembly can likewise bring to God the words of people whose responsibility for its condition differs greatly. Some in Isaiah's prayer may speak with the assurance of the saints. The sinner who would hold God accountable for the hardness of his own heart can also be heard in it. What brings them together at the close is a claim they can make about God whatever their conduct has been. They are all the work of his hands. Even someone who can find no excuse can appeal from his own plight to the maker who is also his Father.\footnote{Jerome, \textit{Commentary on Isaiah}, XVII, on Is 63:16--19 and 64:7--8, Latin part IV, pp.~149--151, 158--159.}

Chrysostom's particular concern is the person in Corinth who needs correction and might stop listening as soon as it begins. Paul recalls what God has done among the Corinthians. That recollection may itself be painful. People who once sought the gifts of God eagerly can hear how generous God was and remember how far their own eagerness has cooled. Chrysostom also considers the possibility that different people are in view. Paul knows that members of this church still have a psalm to offer or teaching to give when they gather. It would be unjust to speak of their services as though all had abandoned their posts. Sometimes too an author says of a whole body what holds of the greater part. On any of these accounts a particular Corinthian can hear the thanksgiving, recognize a real gift at work among his neighbors and still need the correction about to be offered to him.\footnote{John Chrysostom, \textit{Homilies on First Corinthians}, Homily 2, §§4--5, on 1 Cor 1:7, 3:1 and 14:26--29.}

The two communities owe God quite different accounts. Isaiah's speakers have neglected even to call on him. In Corinth there is no shortage of people who can speak. Someone who has hardly begun to pray can join the assembly in asking God for help. Someone who can already teach others to pray has to consider what will happen to his learners if his teaching wins them to a party. By thanking God for the gift, Paul can address both the people still being helped by it and the teacher who may be abusing it.

\subsection{What it means to be found doing right}

In the U.S. proclamation Isaiah's people wish that God would find them doing right. The wish is followed closely by an admission that even the things they count to their credit are stained. In Jerome's Latin text, also reflected in the Challoner Bible at Isaiah 64:5, God has already met someone who does justice. Moreover, following their confession the people express confidence that they will be saved. Those are words Jerome has to account for. The sinner has lived with his sin for a long time. He might cite God's forbearance as evidence that God permitted it. Jerome imagines a physician keeping a sharp instrument from the sight of a frightened patient. A patient who took the delay for indifference would have misunderstood the one who meant to heal him. The ground for confidence Jerome finds in the Latin assurance is the mercy of the one who waited. In the U.S. passage the people pass from the confession of sin directly to their uncleanness. They have particular reason to desire what the Collect asks, the will to go out and meet Christ with deeds he can judge to be good.\footnote{Jerome, \textit{Commentary on Isaiah}, XVII, pp.~149--151, 155. The U.S. Is 64:4 and Challoner Is 64:5 have different verse numbers and different clauses here.}

Thomas Aquinas explains how much might have to change for a deed of justice to be sound. A good undertaking can be carried out negligently. It can be undertaken for a bad reason. Pride can enter when something has been carried out well. Suppose the responsibility is providing food for somebody in need. It matters very much to that person whether the food ever arrives. It matters to the giver whether his first thought in providing it was to be seen in a good light. Even someone who has given for the best reason can begin to regard the work as a distinction he has earned for himself. The petition for a willing heart concerns the entire time a good deed takes, including the moment one first looks back on it.\footnote{Thomas Aquinas, \textit{Exposition on Isaiah}, chapter 64, §85024.}

What if a servant has become accustomed to doing none of the things he was entrusted to do? The vineyard prayer in Psalm 80 permits him a very particular request. In the psalm as Augustine knew it, the worshippers first ask God to turn them. When they speak of the vineyard, they ask him to turn back toward what he planted. Why address both requests to God? The planter was responsible for the vine's beginning and can bring his work to completion. He can give life to people in whom the will to call on him has almost died. Their willingness to ask for that help may already be evidence of it. Peter believed he could be loyal even if all the other disciples failed. In the garden he discovers what a friend may still have to ask of God in the very hour when he wishes to be faithful.\footnote{Augustine, \textit{Exposition on Psalm} 79, §§3, 9, 11. Mark 14:29--42.}

Jerome does something surprising when he reaches Isaiah's admission that God has hidden his face. He begins to pray with his readers. They can acknowledge their own sin, hear God's promise to blot it out, and then ask with the psalmist to see his face. A person who was listening to Isaiah in order to decide who was at fault now has words in which to confess a fault of his own. In the Sunday response the person beside him sings those words too. Nothing requires either of them to know what the other has come to confess.\footnote{Jerome, \textit{Commentary on Isaiah}, XVII, p.~158, on Is 64:7, citing Ps 79:4.}

What if the sight for which they ask has already been given? Augustine recalls people looking at Jesus on the cross. Anyone standing there could see his weakness. Recognizing the Lord in the man who hung there would have required seeing more than the people who mocked him saw. The prayer for his face to shine can be made by someone who wants to recognize the one who has already drawn near. The bishops make a related point about the Collect. On her way to church a worshipper is setting out to meet Christ. She will hear his word and receive Communion. She may also meet someone who needs her help. She could arrive at the church door in good time and still have passed by someone in whom Christ came to meet her.\footnote{Augustine, \textit{Exposition on Psalm} 79, §3. USCCB, \textit{Commentary on the Proper Prayers of Advent}, First Sunday, Collect.}

\subsection{Which gifts belong to a servant?}

A church that knows its neighbors' gifts intimately can be especially vulnerable to envy. A speaker who gets a hearing can make another member feel that she has little to contribute. Thomas Aquinas draws on Augustine when explaining what Paul says in 1 Corinthians 1:5 about speech and knowledge. A member who loves the unity of the Church already has a share in her neighbor's ability. She can be glad that a person who needs instruction will be well taught, even if she could never teach that person herself. Envy deprives her of that good for as long as it makes another's excellence an injury in her eyes. On verse 7 Thomas asks how all the gifts can be present in one church. God can give different gifts to different people. Where Chrysostom considers what has become of people once richly endowed, Thomas explains how Paul can speak in the present of an abundance to which many members contribute.\footnote{Thomas Aquinas, \textit{Commentary on First Corinthians}, chapter 1, lecture 1, §86234, on vv.~5 and 7. Chrysostom, \textit{Homilies on First Corinthians}, Homily 2, §§4--5.}

The vine about which the people sing also has a past that forbids anyone to claim what he has received as exclusively his. Augustine remembers God's promise to Abraham that all nations would be blessed in his seed. He thinks of the Gentiles through Paul's comparison with wild olive branches grafted onto an older root. They live because they have become part of something already planted. Augustine understands Christ as the foundation on which God perfects his planting. Bellarmine recalls the actual people through whom the old vineyard was renewed. The apostles were Israelites. So were thousands who came to faith through Peter's preaching. Gentiles entered that company by being grafted onto Israel's tree. When Bellarmine asks God to give a better vineyard in the old one's place, these are the people among whom he understands God to have reformed the vineyard. The Son of Man has charge of it as its chief husbandman. Both interpreters hear Israel's appeal for attention to God's own work. In one, the question is how the promise to Abraham reaches all the nations. In the other it is who tends the vineyard in which both Israelites and Gentiles now grow.\footnote{Augustine, \textit{Exposition on Psalm} 79, §9, with Gen 22:18 and 1 Cor 3:11. Robert Bellarmine, \textit{A Commentary on the Book of Psalms}, Psalm 79, explanations 14--15, trans. John O'Sullivan. Both commentators draw on Romans 11.}

\subsection{Who is keeping the door?}

Bede takes the householder's journey to mean the bodily Ascension of Christ. A master who sets out on the road would ordinarily be beyond his servants' reach. Christ can be with his Father in the flesh he took from us and still give his servants the Spirit by whom they do their work. Chrysostom reasons from Paul's choice of a word for what lies ahead of the Corinthians. What they expect is a revelation of Christ. Someone can be revealed at a certain time because, already present, he was hidden from those who will then see him. Paul even promises these divided people the help they need to stand firm until then. The promise itself should alert them to their unsteadiness. To a person unsure whether the Lord can still have dealings with such a church, his fidelity is a reason to undertake the correction Paul asks of them.\footnote{Bede, \textit{Commentary on Mark}, on Mark 13:33--37, lines 403--405 of the Latin transcription. Chrysostom, \textit{Homilies on First Corinthians}, Homily 2, §§6--8.}

The servants in the Gospel have distinguishable duties. Bede sees in the porter a pastor whose negligence can expose everyone inside to harm. But Jesus' last word extends the charge to everyone. Each person has a door at which a thought can present itself and obtain consent. A gifted speaker who faithfully ministers to those entrusted to him still has that door of his own to guard. He may receive the admiration of his hearers when he speaks and then be tempted to cultivate a following. Nobody exercising oversight over the community can watch in his place at the moment he begins to entertain that desire. The person with no pastoral office has the same responsibility for the envy that can develop when she hears him praised.

Augustine asks about Christians who will have died long before the final day. If Jesus' warning was of consequence only to the generation that would see his return, most of the Church would have no need of it. In his letter to Hesychius, Augustine recalls the possible hours of arrival in Mark. Every generation has its own period in which someone can take up a charge and carry it out or leave it undone. Death brings that person's opportunity to act to an end. At the judgment still to come, it will matter how he lived and died. Bede asks whose office it is to keep watch tonight. Augustine asks which generations will have to answer for their watch. A person who heard Jesus speak on the mountain had as much cause to attend to him as someone alive on the last day.\footnote{Bede, \textit{Commentary on Mark}, on Mark 13:34--37. Augustine, \textit{Letter} 199 to Hesychius, I.2--3.}

In the garden there is a chance to see what each account of watching asks of the same man. Peter, James and John have been with Jesus on the mountain when he told them of the master's journey. On this later night he takes the three apart from the others. He tells them of his distress and asks them to stay awake. Peter at least has decided what kind of friend he means to be. Only a little earlier he declared himself willing to die with Jesus. Yet when Jesus comes back to him, even the hour for which he wanted company has been too long. The man who spoke so confidently now has good reason to make Israel's petition his own. He needs help to do something he sincerely wished to do. This time Jesus expressly tells his friends to pray against temptation.\footnote{Mark 13:3--4, 33--37 and 14:29--42.}

Bede's account of the porter also invites us to consider what that hour was like for the one who asked Peter to stay. Jesus returned to find his friends sleeping while he was in distress. Peter's willingness to undertake some greater service in the future had given Jesus no company in the hour he had asked him to share. A person hearing the Sunday Gospel might wonder how to spend a night waiting for someone whose hour is unknown. In the garden the friends have someone already there with them. They know enough about his need to keep him company while he prays. Someone who has another person's welfare in his hands can begin his watch that way.

Who will the servants see when the door opens? At Mass the congregation asks in the Collect to stand at Christ's right hand. Cyril of Jerusalem, teaching the Creed, recalls how Jesus once stood before those who judged him and kept silence. At his return he will be the one who judges. For Peter the man before whom his conduct lay open was someone who had sought his friendship in the garden. For a Corinthian who has been quarreling over a preacher, Paul has a question about the same Lord. Who was crucified for him? The preacher whose side he took owes his life to the same Lord who was crucified for the people in the rival party. A Corinthian hoping to meet Christ could begin by making peace with someone to whom Christ had also given a place in his Church.\footnote{Cyril of Jerusalem, \textit{Catechetical Lecture} XV, §§1--3. First Sunday of Advent, Collect and Preface I. 1 Corinthians 1:9--13.}

\subsection{What grows when God gives?}

At Communion, where the congregation hears of the earth yielding fruit, Augustine and Bellarmine differ over what the fruit is. The preceding verse speaks of truth springing from the earth. Augustine can understand that truth as Christ born of a woman or as the truth a sinner speaks when at last he confesses. Augustine tells his hearers that their confession has opened their hearts as a plough opens the ground. He has sown the word there and they have listened. He asks God to send the rain that will make the seed grow after they go home. At the verse about fruit, the word Augustine has for God's gift suggests sweetness. The thief in his example once enjoyed taking what someone else had and he lacked. Now he has something another person needs. He takes pleasure in giving it away, even though he has less afterward. Augustine understands the pleasure of doing justice as something God can give.\footnote{Augustine, \textit{Exposition on Psalm} 84, §§10--12. The preceding truth verse and the appointed fruit verse are Challoner Ps 84:12 and 13 respectively.}

Bellarmine pictures a flower coming up in open ground. Nobody has planted it there. Sun and rain have reached it from above. He identifies the land with Mary, the gift from above with the work of the Holy Spirit, and the fruit with Christ himself. The one the people asked to see when they sang of salvation before the Gospel has already come to them in human flesh. At Communion the two explanations give a communicant different things to consider. If, with Bellarmine, she sees Christ in the fruit, the hands she extends receive the very person for whose coming the Church gives thanks in the Preface. If, with Augustine, she asks about the sweetness of doing right, the question concerns what will become possible in her dealings with someone else once she has gone home. Augustine's account of the preceding verse about truth leaves room for the Incarnation even as he explains the fruit verse by a change in human desire.\footnote{Bellarmine, \textit{Commentary on the Book of Psalms}, Psalm 84, explanations 7 and 12. Augustine, \textit{Exposition on Psalm} 84, §§6--7, 10--12.}

The bread offered at Mass brings those questions close to one another. Human labor was needed to make it. Yet the people who offer it depend both on God's gifts for the means of making bread and on Christ for the gift of his Body. After Communion a person who can afford to feed someone will have another occasion to use something that perishes. It is quite possible to see the need clearly and begrudge what it would cost to meet it. Such a person has good reason to ask God for the will to part with what she has. Someone who does enjoy giving still needs to consider how the food will reach the person who needs it. The master in Jesus' story had people inside his house whose welfare depended on someone actually keeping his post.\footnote{USCCB, \textit{Commentary on the Proper Prayers of Advent}, First Sunday, Prayer over the Offerings and Prayer after Communion.}
